\documentclass[11pt,a4paper]{article}

\usepackage[T1]{fontenc}
\usepackage[utf8]{inputenc}
\usepackage{lmodern}
\usepackage{microtype}
\usepackage{amsmath,amssymb,amsthm,mathtools}
\usepackage{booktabs,array}
\usepackage[margin=2.5cm]{geometry}
\usepackage{enumitem}
\usepackage{xcolor}
\usepackage[colorlinks=true,linkcolor=blue!60!black,citecolor=blue!60!black,urlcolor=blue!60!black]{hyperref}

\newcommand{\R}{\mathbb{R}}
\newcommand{\C}{\mathbb{C}}
\newcommand{\D}{\mathbb{D}}
\newcommand{\N}{\mathbb{N}}
\newcommand{\E}{\mathbb{E}}
\newcommand{\Prob}{\mathbb{P}}
\newcommand{\G}{\mathcal{G}}
\newcommand{\sgn}{\operatorname{sgn}}
\newcommand{\RePart}{\operatorname{Re}}
\newcommand{\ImPart}{\operatorname{Im}}
\newcommand{\adj}{\operatorname{adj}}
\newcommand{\diag}{\operatorname{diag}}
\newcommand{\erf}{\operatorname{erf}}
\newcommand{\Lone}[1]{\lVert #1\rVert_1}
\urlstyle{tt}
\renewcommand{\UrlFont}{\ttfamily\small}
\makeatletter
\g@addto@macro\UrlBreaks{\do\_\do\.}
\makeatother
\newcommand{\code}[1]{\nolinkurl{#1}}

\newtheorem{theorem}{Theorem}
\newtheorem{proposition}{Proposition}[section]
\newtheorem{lemma}[proposition]{Lemma}
\newtheorem{corollary}[proposition]{Corollary}
\theoremstyle{remark}
\newtheorem{remark}[proposition]{Remark}

\numberwithin{equation}{section}
\pagestyle{plain}
\emergencystretch=1.5em

\title{From Quadratic Polynomials to Arbitrary Order:\\
Sharp $L^1$ Interpolation Inequalities and the Center Maximum}
\author{Lijails, GPT 6.2 Astra, GPT 6.2 Sol, Gemini 4 Neon, and Gemini 4 Hydrogen}
\date{}

\begin{document}
\maketitle

\begin{abstract}
For every integer $n \ge 1$ and every compact interval $[a,b]$ of length $L = b - a > 0$, we determine the sharp pair of constants in the $L^1$ Sobolev interpolation inequality
\[
    \lVert f^{(n)}\rVert_{L^1([a,b])}
    \le 4^n n!\,L^{-n}\lVert f\rVert_{L^1([a,b])}
      + \frac{L}{2}B_n\lVert f^{(n+1)}\rVert_{L^1([a,b])}
    \qquad \bigl(f \in W^{n+1,1}([a,b])\bigr),
\]
where the first constant is minimal across all valid inequalities and $B_n$ is minimal once the first constant is fixed. On the reference interval $I=[-1,1]$, the first best constant $A_n = 2^n n!$ is governed by the second-kind Chebyshev polynomial $U_n$, and a moving-node operator $T_n$ driven by the eigenvalues of a symmetric tridiagonal Jacobi matrix $J_n(x)$ reduces the second best constant to $B_n = \max_{t \in [-1,1]} E_n(t)$ for an explicit even error profile $E_n$.

Our main analytic contribution is an unconditional proof for every $n \ge 1$ that $E_n$ is strictly decreasing on $[0,1]$, hence attains its unique global maximum at the center $t=0$ with exact value
\[
    B_n = E_n(0)
    = 1 - \frac{2^n}{n+1}
      + \frac{2^{n+1}}{n+1}\sum_{k=1}^{\lfloor(n+1)/2\rfloor}(-1)^{k-1}\cos^{n+1}\frac{k\pi}{n+2}.
\]
After clearing the exterior interval $[1/2,1)$ directly from node-speed monotonicity, we express $-E_n'(t)$ on $(0,1/2]$ through a complex Gaussian quadratic form whose characteristic function collapses in Joukowsky coordinates to a closed rational function. Deforming the Fourier inversion contour onto the vertical chord $u = t + iv$ eliminates oscillation, exposes an exponential $(3/4)^n$ versus $(1/4)^n$ separation between the half-line main term and the finite-boundary reflection tail, and resolves all orders $n \ge 4$ via closed-form Beta bounds; orders $n=1,2$ follow from elementary algebra and $n=3$ from a five-subinterval rigorous bound. As a byproduct of the unique center maximum, equality in $C^{n+1}([a,b])$ holds if and only if $f$ is a scalar multiple of the rescaled Chebyshev polynomial $U_n$.

The entire proof of Theorem~\ref{thm:full} has been formally verified in Lean~4 / Mathlib with zero \code{sorry}, zero \code{native_decide}, and only the standard Lean axioms (\code{propext}, \code{Classical.choice}, \code{Quot.sound}).
\end{abstract}

\section{Introduction: From Second Order to Arbitrary Order}
\label{sec:intro}

Start on the reference interval $I = [-1,1]$, writing $\Lone{g} := \int_{-1}^{1} |g(t)|\,dt$ for the $L^1(I)$ norm. In the second-order case $n = 2$, one encounters the sharp inequality
\begin{equation}
    \Lone{F''} \le 8\Lone{F} + \frac{2\sqrt{2}-1}{3}\Lone{F'''}
    \qquad \bigl(F \in W^{3,1}(I)\bigr),
    \label{eq:n2}
\end{equation}
accompanied by two striking phenomena. First, the leading constant $8$ is forced already on quadratic polynomials by the monic Chebyshev polynomial $x^2 - \frac14$. Second, once the constant $8$ is locked in, constructing a three-point moving-node approximation $T$ to $F''$ produces an error profile $E(t)$ whose unique global maximum on $[-1,1]$ sits right at the midpoint $t = 0$, yielding the exact second constant $\frac{2\sqrt{2}-1}{3}$.

A natural question immediately arises. Replace $F''$ by an arbitrary derivative $F^{(n)}$ ($n \ge 1$), and seek first the smallest possible constant $A_n$, and then---with $A_n$ fixed---the smallest possible constant $B_n$ such that
\begin{equation}
    \Lone{F^{(n)}} \le A_n\Lone{F} + B_n\Lone{F^{(n+1)}}
    \qquad \bigl(F \in W^{n+1,1}(I)\bigr).
    \label{eq:main}
\end{equation}
What are the exact optimal constants $A_n$ and $B_n$? Above all, does the worst-case jump of $F^{(n)}$ remain at the midpoint $t = 0$ for \emph{every} order $n \ge 1$?

We answer all of these questions affirmatively and unconditionally.

\begin{theorem}[Sharp $L^1$ interpolation inequality and center maximum]
    \label{thm:full}
    Let $n \ge 1$ be an integer, and let $E_n : [-1,1] \to \R$ be the moving-node error profile defined in~\eqref{eq:En-def} and explicitly evaluated in~\eqref{eq:En}.
    \begin{enumerate}[label=\textup{(\roman*)}]
        \item \textbf{Strict monotonicity and unique center maximum.}
        For every $t \in (0,1)$, $E_n'(t) < 0$. Consequently, $E_n$ is strictly decreasing on $[0,1]$ and attains its unique global maximum on $[-1,1]$ at $t = 0$, with exact value
        \begin{equation}
            B_n := E_n(0)
            = 1 - \frac{2^n}{n+1}
              + \frac{2^{n+1}}{n+1}\sum_{k=1}^{\lfloor(n+1)/2\rfloor}(-1)^{k-1}\cos^{n+1}\frac{k\pi}{n+2}.
            \label{eq:bn-thm}
        \end{equation}
        \item \textbf{Sharp $W^{n+1,1}$ inequality on arbitrary intervals.}
        For every compact interval $[a,b]$ of length $L = b - a > 0$ and every $f \in W^{n+1,1}([a,b])$,
        \begin{equation}
            \lVert f^{(n)}\rVert_{L^1([a,b])}
            \le 4^n n!\,L^{-n}\lVert f\rVert_{L^1([a,b])}
              + \frac{L}{2}B_n\lVert f^{(n+1)}\rVert_{L^1([a,b])}.
            \label{eq:full-ab}
        \end{equation}
        \item \textbf{Optimality of both constants.}
        Even when restricted to $f \in C^\infty([a,b])$, the first constant $4^n n!\,L^{-n}$ cannot be replaced by any smaller constant (regardless of the second constant); once the first constant is fixed at $4^n n!\,L^{-n}$, the second constant $\frac{L}{2}B_n$ cannot be replaced by any smaller constant.
        \item \textbf{Exact $C^{n+1}$ equality characterization.}
        For $f \in C^{n+1}([a,b])$, equality holds in~\eqref{eq:full-ab} if and only if there exists $\lambda \in \R$ such that
        \[
            f(x) = \lambda\, U_n\!\left(\frac{2x - a - b}{b - a}\right)
            \qquad \text{for all } x \in [a,b],
        \]
        where $U_n$ denotes the $n$-th Chebyshev polynomial of the second kind.
    \end{enumerate}
\end{theorem}

The proof follows the same transparent rhythm as the second-order prototype, punctuated by one decisive complex-analytic turn:
\begin{itemize}[leftmargin=1.8em]
    \item In Section~\ref{sec:poly}, setting $F^{(n+1)} \equiv 0$ isolates the polynomial inequality, identifies the extremal polynomial $Q_n = 2^{-n}U_n$, and fixes $A_n = 2^n n!$.
    \item In Section~\ref{sec:matrix}, the sign intervals of $U_n$ motivate an $(n+1)$-point moving-node operator $T_n$, whose node trajectories are realized globally as the eigenvalues of a symmetric tridiagonal Jacobi matrix $J_n(x)$.
    \item In Section~\ref{sec:sharp}, Taylor's integral remainder expresses $F^{(n)} - T_n F$ against the error profile $E_n(t)$, while a sign-aligned spline $G_t = h_t - P_t$ and smooth mollification prove that $B_n = \max_{t \in I} E_n(t)$ is sharp.
    \item In Section~\ref{sec:setup}, direct differentiation immediately disposes of the exterior region $t \in [1/2,1)$. On the interior $t \in (0,1/2]$, where the raw alternating sum suffers exponential cancellation, rewriting $-E_n'(t)$ as a complex Gaussian quadratic form splits it into a strictly positive symmetry term plus $\frac{n}{2}\Delta_n(t)$.
    \item In Section~\ref{sec:chord}, the characteristic function of the quadratic form simplifies in Joukowsky coordinates to a closed rational function, allowing Cauchy contour deformation onto the vertical chord $u = t + iv$. Oscillation vanishes on this chord.
    \item In Sections~\ref{sec:n4} and~\ref{sec:finish}, separating real and imaginary parts preserves the vanishing factor $t$ as $t \to 0^+$, revealing a $(3/4)^n$ main term against a $(1/4)^n$ reflection tail that settles all $n \ge 4$ uniformly; orders $n = 1, 2$ are elementary, $n = 3$ is certified on five compact subintervals, and Theorem~\ref{thm:full} follows.
    \item Section~\ref{sec:lean} documents the complete Lean~4 formalization and gives an exact module-by-module dictionary, while Section~\ref{sec:open} reflects on the proof and records open questions. Appendices~\ref{app:extensions} and~\ref{app:tech} present expository structural complements (generalized Gauss quadrature, a cancellation-free integral and $1/\sqrt{\pi n}$ asymptotics for $B_n$, and the Gaussian scaling profile) alongside the technical polynomial and $n = 3$ certification details.
\end{itemize}

\section{Without the \texorpdfstring{$(n+1)$}{(n+1)}-st Derivative: Chebyshev Polynomials}
\label{sec:poly}

Set $F^{(n+1)} \equiv 0$, so that $F = P$ is a polynomial of degree at most $n$. Inequality~\eqref{eq:main} reduces to
\begin{equation}
    \Lone{P^{(n)}} \le A_n \Lone{P} \qquad (\deg P \le n).
    \label{eq:poly-ineq}
\end{equation}
The left-hand side sees only the leading coefficient $a_n$ of $P(y) = a_n y^n + \dots + a_0$, namely $\Lone{P^{(n)}} = 2n!|a_n|$, whereas the right-hand side mixes in all lower-degree terms. Just as in the quadratic case, we seek a $\pm 1$-valued sign function $s_n(y)$ on $[-1,1]$ that annihilates every monomial of degree strictly less than $n$:
\begin{equation}
    \int_{-1}^{1} s_n(y)\,y^k\,dy = 0 \qquad (0 \le k < n).
    \label{eq:sn-ortho}
\end{equation}

For $n = 2$, $s_2$ was the sign of $x^2 - \frac14$. For general $n \ge 1$, the answer is the sign of the \emph{Chebyshev polynomial of the second kind} $U_n$, characterized on $[-1,1]$ by $U_n(\cos\theta) = \frac{\sin((n+1)\theta)}{\sin\theta}$. The reason is immediate under the substitution $y = \cos\theta$:
\[
    \int_{-1}^{1} \sgn\bigl(U_n(y)\bigr)\,y^k\,dy
    = \int_0^\pi \sgn\sin\bigl((n+1)\theta\bigr)\cdot \cos^k\theta\,\sin\theta\,d\theta.
\]
The square wave $\sgn\sin((n+1)\theta)$ contains only odd multiples of the frequency $n+1$ in its sine expansion (equivalently, shifting $\theta \mapsto \theta + \frac{\pi}{n+1}$ negates it, or testing directly against $\sin(m\theta)$ for $1 \le m \le n$ gives zero by elementary cancellation across the $n+1$ equal subintervals of $[0,\pi]$). Meanwhile, $\cos^k\theta\,\sin\theta$ is a linear combination of $\sin(m\theta)$ with $1 \le m \le k+1$. Whenever $k < n$, the frequencies never meet, so~\eqref{eq:sn-ortho} holds identically.

Define the monic rescaling $Q_n := 2^{-n} U_n$ and its sign function $s_n := \sgn Q_n$. The same trigonometric substitution computes the exact $L^1$ norm:
\[
    \Lone{Q_n}
    = 2^{-n}\int_0^\pi \bigl|\sin\bigl((n+1)\theta\bigr)\bigr|\,d\theta
    = 2^{-n}(n+1)\int_0^{\pi/(n+1)} \sin\bigl((n+1)\theta\bigr)\,d\theta
    = 2^{1-n}.
\]

\begin{proposition}[First best constant]
    \label{prop:An}
    For every integer $n \ge 1$ and every polynomial $P$ of degree at most $n$,
    \[
        \Lone{P^{(n)}} \le 2^n n!\,\Lone{P},
    \]
    with equality if and only if $P = \lambda U_n$ for some $\lambda \in \R$. In particular, any valid constant $A$ in~\eqref{eq:main} satisfies $A \ge A_n := 2^n n!$.
\end{proposition}

\begin{proof}
Write $P(y) = a_n Q_n(y) + R(y)$ with $\deg R < n$. By~\eqref{eq:sn-ortho}, $\int_{-1}^{1} s_n(y)R(y)\,dy = 0$, while $\int_{-1}^{1} s_n(y)Q_n(y)\,dy = \Lone{Q_n} = 2^{1-n}$. Hence
\[
    \Lone{P^{(n)}}
    = 2n!|a_n|
    = 2^n n!\left|\int_{-1}^{1} s_n(y)P(y)\,dy\right|
    \le 2^n n!\int_{-1}^{1} |P(y)|\,dy
    = 2^n n!\,\Lone{P}.
\]
Taking $P = \lambda U_n$ makes $s_n(y)P(y)$ have constant sign almost everywhere, so equality holds. Conversely, if equality holds for a nonzero polynomial $P$ of degree at most $n$, then $|P(y)| - \sigma s_n(y)P(y) = 0$ almost everywhere for some $\sigma \in \{-1, 1\}$. Since $P$ is continuous and $s_n$ changes sign at each of the $n$ interior zeros $r_1 < \dots < r_n$ of $U_n$, $P$ must vanish at all $n$ zeros of $U_n$. Because $\deg P \le n$, it follows that $P = \lambda U_n$ for some $\lambda \in \R$. Finally, evaluating~\eqref{eq:main} on $F = Q_n$ (for which $Q_n^{(n+1)} \equiv 0$) yields $A \ge \Lone{Q_n^{(n)}}/\Lone{Q_n} = 2^n n!$.
\end{proof}

\section{Moving Nodes and Jacobi Matrices: From Three Points to \texorpdfstring{$n+1$}{n+1} Points}
\label{sec:matrix}

Armed with the extremal polynomial $Q_n = 2^{-n}U_n$, we follow the second-order blueprint. List the $n$ zeros of $U_n$ together with the endpoints $\pm 1$ in increasing order:
\begin{equation}
    r_j := -\cos\frac{j\pi}{n+1} \qquad (j = 0, 1, \dots, n+1).
    \label{eq:rj}
\end{equation}
They partition $[-1,1]$ into $n+1$ subintervals $[r_j, r_{j+1}]$ on which $s_n$ takes the alternating signs $\varepsilon_j := (-1)^{n-j}$. Our goal is to construct $n+1$ strictly increasing smooth bijections $u_j : [-1,1] \to [r_j, r_{j+1}]$ ($j = 0, \dots, n$) and define the \emph{moving-node operator}
\begin{equation}
    T_n F(x) := 2^n n!\sum_{j=0}^{n} \varepsilon_j\, u_j'(x)\, F\bigl(u_j(x)\bigr).
    \label{eq:Tn}
\end{equation}
Regardless of the specific choice of strictly increasing bijections $u_j : [-1,1] \to [r_j, r_{j+1}]$, the triangle inequality and piecewise change of variables $y = u_j(x)$ immediately give the sharp $L^1$ contraction bound
\begin{equation}
    \Lone{T_n F}
    \le 2^n n!\sum_{j=0}^{n}\int_{-1}^{1} u_j'(x)\bigl|F(u_j(x))\bigr|\,dx
    = 2^n n!\sum_{j=0}^{n}\int_{r_j}^{r_{j+1}} |F(y)|\,dy
    = A_n \Lone{F}.
    \label{eq:Tn-bound}
\end{equation}

To make $T_n$ reproduce $F^{(n)}$ on polynomials of degree at most $n$, we demand $T_n P = P^{(n)}$ for all $\deg P \le n$, which is equivalent to the moment system
\[
    \sum_{j=0}^{n} \varepsilon_j\, u_j'(x)\, u_j(x)^k = 0 \quad (0 \le k < n),
    \qquad
    \sum_{j=0}^{n} \varepsilon_j\, u_j'(x)\, u_j(x)^n = 2^{-n}.
\]
Because the Vandermonde system on $n+1$ distinct ordered nodes $u_0(x) < \dots < u_n(x)$ has a unique solution---namely the barycentric weights of Lagrange interpolation---solving for the node velocities yields
\begin{equation}
    u_j'(x) = \frac{2^{-n}}{\prod_{i \neq j} |u_j(x) - u_i(x)|},
    \qquad
    T_n F(x) = n!\, F\bigl[u_0(x), \dots, u_n(x)\bigr],
    \label{eq:ode}
\end{equation}
where $F[y_0, \dots, y_n] = \sum_{j=0}^n \frac{F(y_j)}{\prod_{i \neq j}(y_j - y_i)}$ denotes the $n$-th divided difference. In short, $T_n F(x)$ is simply $n!$ times the $n$-th divided difference of $F$ along the moving nodes $u_0(x), \dots, u_n(x)$.

In the second-order case $n = 2$, one could integrate the three ODEs directly and solve a quadratic equation. For arbitrary $n$, integrating~\eqref{eq:ode} by hand looks daunting---until a remarkably clean linear-algebraic object resolves everything at once. Consider the $(n+1) \times (n+1)$ real symmetric tridiagonal \emph{Jacobi matrix}
\begin{equation}
    J_n(x) := \frac12
    \begin{pmatrix}
        x & 1 & & & \\
        1 & 0 & 1 & & \\
        & 1 & \ddots & \ddots & \\
        & & \ddots & 0 & 1 \\
        & & & 1 & x
    \end{pmatrix}.
    \label{eq:Jn}
\end{equation}
As an immediate sanity check at $n = 2$, the $3 \times 3$ matrix $\frac12\bigl(\begin{smallmatrix} x & 1 & 0 \\ 1 & 0 & 1 \\ 0 & 1 & x \end{smallmatrix}\bigr)$ has eigenvalues $\frac{x - \sqrt{x^2+8}}{4}$, $\frac{x}{2}$, and $\frac{x + \sqrt{x^2+8}}{4}$---word for word the three moving nodes of~\eqref{eq:n2}.

\begin{proposition}[Jacobi realization of the moving nodes]
    \label{prop:jacobi}
    For every integer $n \ge 1$ and every $x \in \R$, let $u_0(x) < u_1(x) < \dots < u_n(x)$ be the eigenvalues of $J_n(x)$ in increasing order. Then:
    \begin{enumerate}[label=\textup{(\roman*)}]
        \item \textbf{Characteristic polynomial and boundary values.}
        The monic characteristic polynomial of $J_n(x)$ is
        \begin{equation}
            p_x(y) := \det\bigl(yI - J_n(x)\bigr)
            = 2^{-n-1}\bigl(U_{n+1}(y) - 2xU_n(y) + x^2 U_{n-1}(y)\bigr).
            \label{eq:charpoly}
        \end{equation}
        In particular, $p_{-1}(y) = 2^{-n}(y+1)U_n(y)$, $p_0(y) = 2^{-n-1}U_{n+1}(y)$, and $p_1(y) = 2^{-n}(y-1)U_n(y)$, so $u_j(-1) = r_j$, $u_j(0) = -\cos\frac{(j+1)\pi}{n+2}$, and $u_j(1) = r_{j+1}$ for every $j = 0, \dots, n$.
        \item \textbf{Automatic solution of the ODE.}
        Each $u_j$ is real-analytic on $\R$, strictly increasing, and satisfies~\eqref{eq:ode}. In particular, $u_j$ maps $[-1,1]$ bijectively onto $[r_j, r_{j+1}]$ and $u_j(x) \in (-1,1)$ for all $x \in (-1,1)$.
        \item \textbf{Reflection symmetry and rational speed formula.}
        For every $x \in \R$ and $j \in \{0, \dots, n\}$,
        \[
            u_j(-x) = -u_{n-j}(x), \qquad u_j'(-x) = u_{n-j}'(x),
        \]
        and the node velocity is a rational function of $x$ and $u_j(x)$:
        \begin{equation}
            u_j'(x) = \frac{2\bigl(1 - u_j(x)^2\bigr)}{L_x\bigl(u_j(x)\bigr)},
            \qquad
            L_x(y) := n + 2 - 2(n+1)xy + nx^2.
            \label{eq:speed}
        \end{equation}
    \end{enumerate}
\end{proposition}

\begin{proof}
(i) Expanding $\det(yI - J_n(x))$ along the first and last rows expresses it in terms of the determinants $d_m(y) = 2^{-m}U_m(y)$ of the unperturbed $m \times m$ tridiagonal Toeplitz blocks with $\frac12$ on the sub- and super-diagonals:
\[
    \det\bigl(yI - J_n(x)\bigr)
    = \Bigl(y - \frac{x}{2}\Bigr)^2 d_{n-1}(y)
      - \frac12\Bigl(y - \frac{x}{2}\Bigr)d_{n-2}(y)
      + \frac{1}{16} d_{n-3}(y)
\]
(with the obvious small-$n$ conventions $d_0 = 1$, $d_{-1} = 0$, $U_{-1} = 0$). Using the Chebyshev recurrence $2yU_m(y) = U_{m+1}(y) + U_{m-1}(y)$ simplifies this immediately to~\eqref{eq:charpoly}. Evaluating at $x = -1, 0, 1$ via $U_{n+1}(y) \pm 2U_n(y) + U_{n-1}(y) = 2(y \pm 1)U_n(y)$ gives the stated factorizations and boundary nodes.

(ii) Because $J_n(x)$ is symmetric tridiagonal with strictly positive off-diagonal entries $\frac12$, all of its eigenvalues are simple, hence real-analytic in $x$. Let $v_j = (v_{j,0}, \dots, v_{j,n})^T$ be a unit eigenvector for $u_j(x)$. Since $J_n(x)$ is centrosymmetric (invariant under reversing coordinate order), $v_j$ is either symmetric or skew-symmetric, so $|v_{j,0}| = |v_{j,n}|$. By first-order eigenvalue perturbation (Hellmann--Feynman),
\[
    u_j'(x) = v_j^T J_n'(x) v_j = \frac12\bigl(v_{j,0}^2 + v_{j,n}^2\bigr) = v_{j,0}^2 > 0.
\]
On the other hand, the rank-one adjugate identity $\adj(u_j I - J_n(x)) = p_x'(u_j)\, v_j v_j^T$ has top-right entry equal to the product of the $n$ off-diagonal entries of $J_n(x)$, namely $2^{-n}$ up to sign. Comparing magnitudes gives $|p_x'(u_j)|\,|v_{j,0} v_{j,n}| = 2^{-n}$, i.e.\ $|p_x'(u_j)|\,u_j'(x) = 2^{-n}$, which is~\eqref{eq:ode}. Since $u_j(-1) = r_j$ and $u_j(1) = r_{j+1}$, strict monotonicity forces $u_j([-1,1]) = [r_j, r_{j+1}]$.

(iii) Conjugating $J_n(x)$ by the product of coordinate reversal and the alternating-sign matrix $S = \diag(1, -1, 1, \dots, (-1)^n)$ negates $J_n(x)$ while flipping $x \mapsto -x$, so $J_n(-x)$ is orthogonally similar to $-J_n(x)$; hence $u_j(-x) = -u_{n-j}(x)$ and $u_j'(-x) = u_{n-j}'(x)$. Finally, differentiating $p_x(u_j(x)) = 0$ gives $u_j'(x) = -\frac{\partial_x p_x(y)}{\partial_y p_x(y)}\big|_{y=u_j(x)}$. Using the standard Chebyshev derivative identity $(1-y^2)U_m'(y) = -m y U_m(y) + (m+1)U_{m-1}(y)$ together with $U_{n+1}(y) - 2xU_n(y) + x^2 U_{n-1}(y) = 0$ cancels $U_n(y)$ and $U_{n-1}(y)$ cleanly to produce~\eqref{eq:speed}.
\end{proof}

Beyond sparing us from integrating an $n$-dimensional nonlinear ODE system, the matrix $J_n(x)$ hides the master key to the center-maximum problem: the exact characteristic polynomial~\eqref{eq:charpoly}.

\section{The Error Profile \texorpdfstring{$E_n(t)$}{E\_n(t)} and Sharpness of the Second Constant}
\label{sec:sharp}

Define the normalized truncated power $h_t(y) := \frac{1}{n!}(y - t)_+^n$, where $z_+ := \max(z, 0)$. For any $F \in W^{n+1,1}(I)$ (meaning $F, F', \dots, F^{(n)}$ are absolutely continuous on $[-1,1]$ with $F^{(n+1)} \in L^1(I)$), Taylor's formula with integral remainder from the left endpoint $-1$ writes
\[
    F(y) = P_F(y) + \int_{-1}^{1} h_t(y)\, F^{(n+1)}(t)\,dt,
\]
where $\deg P_F \le n$ and $P_F^{(n)} \equiv F^{(n)}(-1)$. Applying $D^n - T_n$ and invoking polynomial exactness $T_n P_F = P_F^{(n)}$ yields the Peano representation
\begin{equation}
    F^{(n)}(x) - T_n F(x)
    = \int_{-1}^{1} \bigl(\mathbf{1}_{\{x > t\}} - T_n h_t(x)\bigr)\, F^{(n+1)}(t)\,dt.
    \label{eq:peano}
\end{equation}
Taking the $L^1$ norm in $x$ and applying Fubini--Tonelli gives
\[
    \Lone{F^{(n)} - T_n F} \le \left(\sup_{t \in I} E_n(t)\right)\Lone{F^{(n+1)}},
\]
where the \emph{error profile} is defined by
\begin{equation}
    E_n(t) := \int_{-1}^{1} \bigl|\mathbf{1}_{\{x > t\}} - T_n h_t(x)\bigr|\,dx
    \qquad (t \in [-1,1]).
    \label{eq:En-def}
\end{equation}

In the second-order case, convexity removed the absolute value inside~\eqref{eq:En-def}. For general $n \ge 1$, the \emph{Hermite--Genocchi formula} for divided differences takes over:
\[
    T_n F(x)
    = n!\, F\bigl[u_0(x), \dots, u_n(x)\bigr]
    = \E\!\left[F^{(n)}\!\left(\sum_{j=0}^{n}\lambda_j u_j(x)\right)\right],
\]
where $(\lambda_0, \dots, \lambda_n)$ is uniformly distributed on the standard $n$-simplex (for $h_t$, whose $n$-th weak derivative is $\mathbf{1}_{\{y > t\}}$, this holds by mollifying $h_t$ and passing to the limit). Because $0 \le \mathbf{1}_{\{y > t\}} \le 1$, it follows immediately that
\begin{equation}
    0 \le T_n h_t(x) \le 1 \qquad \text{for all } x, t \in [-1,1].
    \label{eq:Tn-ht-bounds}
\end{equation}
Thus the integrand in~\eqref{eq:En-def} is $T_n h_t(x)$ for $x \le t$ and $1 - T_n h_t(x)$ for $x > t$. Splitting the integral at $x = t$ and changing variables $y = u_j(x)$ term by term in~\eqref{eq:Tn}, all interior endpoint evaluations at $r_1, \dots, r_n$ cancel in alternating pairs, leaving the closed formula
\begin{equation}
    E_n(t)
    = 1 - t - \frac{2^n}{n+1}(1 - t)^{n+1}
      + \frac{2^{n+1}}{n+1}\sum_{j=0}^{n} \varepsilon_j \bigl(u_j(t) - t\bigr)_+^{n+1}.
    \label{eq:En}
\end{equation}
From the reflection symmetry $u_j(-x) = -u_{n-j}(x)$ and $\varepsilon_{n-j} = (-1)^n \varepsilon_j$, one checks directly that $T_n h_t(x) + T_n h_{-t}(-x) = 1$, so $E_n$ is even:
\[
    E_n(-t) = E_n(t) \qquad (t \in [-1,1]).
\]
Evaluating~\eqref{eq:En} at the center $t = 0$, where $u_j(0) = -\cos\frac{(j+1)\pi}{n+2}$ by Proposition~\ref{prop:jacobi}(i), only the strictly positive nodes ($u_j(0) > 0$) survive, yielding the exact center value
\begin{equation}
    b_n := E_n(0)
    = 1 - \frac{2^n}{n+1}
      + \frac{2^{n+1}}{n+1}\sum_{k=1}^{\lfloor(n+1)/2\rfloor}(-1)^{k-1}\cos^{n+1}\frac{k\pi}{n+2}.
    \label{eq:bn}
\end{equation}

Before maximizing $E_n$, we verify that no slack was introduced by passing through $T_n$: the supremum of $E_n$ is genuinely the sharp second constant.

\begin{proposition}[Sharpness of the second constant]
    \label{prop:sharp}
    Let $n \ge 1$. If a constant $B \in \R$ satisfies
    \[
        \Lone{F^{(n)}} \le A_n\Lone{F} + B\Lone{F^{(n+1)}}
        \qquad \text{for all } F \in C^\infty([-1,1]),
    \]
    with $A_n = 2^n n!$, then $B \ge E_n(t)$ for every $t \in [-1,1]$. In particular, with $A_n$ fixed, the optimal second constant in~\eqref{eq:main} is strictly equal to $B_n = \max_{t \in [-1,1]} E_n(t)$.
\end{proposition}

\begin{proof}
Fix any interior point $t \in (-1,1)$, let $P_t$ be the unique polynomial of degree at most $n$ interpolating $h_t$ at the $n+1$ distinct nodes $u_0(t) < \dots < u_n(t)$, and set $G_t := h_t - P_t$. By the Newton interpolation remainder formula,
\[
    G_t(y) = h_t\bigl[u_0(t), \dots, u_n(t), y\bigr]\prod_{j=0}^{n}\bigl(y - u_j(t)\bigr).
\]
Since the $(n+1)$-st distributional derivative of $h_t$ is the positive Dirac mass $\delta_t$, the $(n+1)$-st divided difference $h_t[u_0(t), \dots, u_n(t), y]$ is nonnegative everywhere (this is verified algebraically in Lean~4 via \code{barycentric_truncatedPower_higher_nonneg}). Consequently, $G_t$ changes sign only at the $n+1$ nodes $u_0(t), \dots, u_n(t)$:
\begin{itemize}[leftmargin=1.8em]
    \item For $x > t$, strict monotonicity of each $u_j$ places $u_j(x) \in (u_j(t), u_{j+1}(t))$ (with $u_{n+1}(t) := 1$), so $\prod_{i=0}^n (u_j(x) - u_i(t))$ has sign $(-1)^{n-j} = \varepsilon_j$. Hence $\varepsilon_j G_t(u_j(x)) \ge 0$ for all $j = 0, \dots, n$.
    \item For $x < t$, $u_j(x) \in (u_{j-1}(t), u_j(t))$ (with $u_{-1}(t) := -1$), so $\varepsilon_j G_t(u_j(x)) \le 0$ for all $j = 0, \dots, n$.
\end{itemize}
Therefore, in the sum defining $T_n G_t(x)$, all $n+1$ terms share the same sign: nonnegative for $x > t$ and nonpositive for $x < t$. Integrating over $[-1,1]$ incurs zero cancellation, so $\Lone{T_n G_t} = A_n \Lone{G_t}$. Moreover, because $P_t$ has degree $\le n$,
\[
    G_t^{(n)}(x) - T_n G_t(x) = \mathbf{1}_{\{x > t\}} - T_n h_t(x)
\]
is likewise nonnegative for $x > t$ and nonpositive for $x < t$ by~\eqref{eq:Tn-ht-bounds}. Hence $T_n G_t(x)$ and $G_t^{(n)}(x) - T_n G_t(x)$ never oppose in sign, which forces exact equality in the triangle inequality $\Lone{G_t^{(n)}} = \Lone{T_n G_t} + \Lone{G_t^{(n)} - T_n G_t}$, i.e.
\[
    \Lone{G_t^{(n)}} - A_n\Lone{G_t} = E_n(t).
\]

To pass from the spline $G_t$ to $C^\infty([-1,1])$, pick a nonnegative bump function $\varphi \in C_c^\infty((-1,1))$ with $\int \varphi = 1$, set $\varphi_\varepsilon(s) := \varepsilon^{-1}\varphi(s/\varepsilon)$ for $0 < \varepsilon < \min(1+t, 1-t)$, and convolve $G_{t,\varepsilon} := G_t * \varphi_\varepsilon$ (extending $G_t$ naturally to $\R$ by $h_t - P_t$). Then $G_{t,\varepsilon} \in C^\infty(\R)$, its $(n+1)$-st derivative is $G_{t,\varepsilon}^{(n+1)}(y) = \varphi_\varepsilon(y - t) \ge 0$ supported inside $(-1,1)$ with $\Lone{G_{t,\varepsilon}^{(n+1)}} = 1$, and as $\varepsilon \to 0^+$ we have $G_{t,\varepsilon} \to G_t$ and $G_{t,\varepsilon}^{(n)} \to G_t^{(n)}$ in $L^1([-1,1])$. Applying the assumed $C^\infty$ inequality to $G_{t,\varepsilon}$ and letting $\varepsilon \to 0^+$ yields $B \ge E_n(t)$ for every $t \in (-1,1)$, and by continuity (or $E_n(\pm 1) = 0$) for all $t \in [-1,1]$.
\end{proof}

The entire sharp constant problem has now distilled to a single assertion:
\begin{quote}\centering\itshape
    For every integer $n \ge 1$, the error profile $E_n(t)$ is strictly decreasing on $[0,1]$, hence attains its unique global maximum on $[-1,1]$ at $t = 0$.
\end{quote}

\section{Attacking Monotonicity: Clearing the Exterior and Translating the Interior}
\label{sec:setup}

Set $H_n(t) := 1 - E_n(t)$, so that our target is $H_n'(t) > 0$ (equivalently $E_n'(t) < 0$) for all $t \in (0,1)$. Differentiating~\eqref{eq:En} (note that $(y - t)_+^{n+1}$ is $C^1$ for every $n \ge 1$, even when a node crosses the diagonal $u_j(t) = t$) yields
\begin{equation}
    E_n'(t) = -1 + 2^n(1 - t)^n
    - 2^{n+1}\sum_{j=0}^{n} \varepsilon_j \bigl(u_j(t) - t\bigr)_+^n \bigl(1 - u_j'(t)\bigr).
    \label{eq:En-prime}
\end{equation}
One glance at the term $2^n(1-t)^n$ reveals a crisp dividing line at $t = \frac12$.

\subsection{The exterior interval \texorpdfstring{$t \in [1/2, 1)$}{t in [1/2, 1)}: a few lines for all orders}
\label{sec:outer}

\begin{proposition}[Exterior strict monotonicity]
    \label{prop:outer}
    For every integer $n \ge 1$ and every $t \in [1/2, 1)$, we have $E_n'(t) < 0$.
\end{proposition}

\begin{proof}
Fix $t \in (0,1)$ and view the rational speed formula~\eqref{eq:speed} as a function of $y \in [t,1]$:
\[
    v(y) := \frac{2(1 - y^2)}{L_t(y)}, \qquad L_t(y) = n + 2 - 2(n+1)ty + nt^2 > 0.
\]
Differentiating gives $v'(y) = -4q(y)/L_t(y)^2$, where $q(y) := yL_t(y) - (n+1)t(1-y^2)$. Direct evaluation shows
\[
    q(t) = t(1 - t^2) > 0,
    \qquad
    q'(y) = n + 2 + nt^2 - 2(n+1)ty \ge q'(1) = (1 - t)\bigl(n(1-t) + 2\bigr) > 0
\]
for all $y \in [t,1]$. Hence $v$ is strictly decreasing on $[t,1]$, and in particular
\[
    v(y) \le v(t) = \frac{2(1 - t^2)}{(n+2)(1 - t^2)} = \frac{2}{n+2} < 1
    \qquad (y \in [t,1]).
\]
It follows that $g(y) := (y - t)_+^n(1 - v(y))$ is nonnegative on $[-1,1]$ and strictly increasing on $[t,1]$. Grouping the alternating sum in~\eqref{eq:En-prime} in pairs starting from the top node $u_n(t)$ (recall $\varepsilon_n = +1$),
\[
    \sum_{j=0}^{n} \varepsilon_j g\bigl(u_j(t)\bigr)
    = \bigl(g(u_n) - g(u_{n-1})\bigr)
    + \bigl(g(u_{n-2}) - g(u_{n-3})\bigr) + \dots \ge 0.
\]
Consequently $E_n'(t) \le -1 + [2(1-t)]^n$, which is already strictly negative for every $t \in (1/2, 1)$. At the boundary $t = 1/2$, direct computation for $n = 1$ gives $E_1'(1/2) = -1/2 < 0$; for $n \ge 2$, the largest node satisfies $u_n(1/2) > u_n(-1) = \cos\frac{\pi}{n+1} \ge \frac12$, so $g(u_n(1/2)) > g(u_{n-1}(1/2))$ and the top pair is strictly positive, again giving $E_n'(1/2) < 0$.
\end{proof}

\subsection{The interior interval \texorpdfstring{$t \in (0, 1/2]$}{t in (0, 1/2]}: Gaussian quadratic-form decomposition}
\label{sec:interior-decomp}

For $t < 1/2$, the term $[2(1-t)]^n$ grows exponentially with $n$, while the alternating sum $\sum_j \varepsilon_j g(u_j(t))$ exhibits massive cancellation between positive and negative terms. Direct termwise estimation of~\eqref{eq:En-prime} is hopeless in this regime.

The way forward is to fold the boundary terms $-1 + 2^n(1-t)^n$ back into the sum so that the entire derivative becomes a combination of divided differences on the $n+3$ extended nodes
\[
    \mathbf{z}_t := \bigl(-1,\, u_0(t),\, \dots,\, u_n(t),\, 1\bigr),
    \qquad
    P_t^*(y) := (y^2 - 1)\,p_t(y).
\]
Comparing~\eqref{eq:ode} and~\eqref{eq:speed}---specifically $|p_t'(u_j)|u_j' = 2^{-n}$ and $u_j' = 2(1-u_j^2)/L_t(u_j)$---together with the endpoint values $P_t^*(\pm 1)' = \pm 2p_t(\pm 1)$ gives the exact structural weight identity
\begin{equation}
    (-1)^{n+1}\delta_{-1} + 2\sum_{j=0}^{n}\varepsilon_j \delta_{u_j(t)} - \delta_1
    = -2^{-n}\sum_{z \in \mathbf{z}_t} \frac{L_t(z)}{(P_t^*)'(z)}\,\delta_z.
    \label{eq:weight-id}
\end{equation}
Expanding the numerator $L_t(y) - 2(1 - y^2)$ around $y = t$ yields
\[
    L_t(y) - 2(1 - y^2) = n(1 - t^2) - 2(n - 1)t(y - t) + 2(y - t)^2,
\]
which collapses~\eqref{eq:En-prime} into three barycentric divided differences on $\mathbf{z}_t$ of $(y-t)_+^{n+2}$, $(y-t)_+^{n+1}$, and $(y-t)_+^{n}$.

Now translate divided differences of truncated powers into a \emph{complex Gaussian quadratic form}. Let $g \sim \mathcal{CN}(0, I_{n+1})$ be a standard complex Gaussian vector in $\C^{n+1}$ (normalized so that $|g_0|^2, \dots, |g_n|^2$ are i.i.d.\ $\mathrm{Exp}(1)$), and let $e_+, e_- \sim \mathrm{Exp}(1)$ be independent standard exponential random variables, also independent of $g$. Define
\begin{equation}
    Q_t^0 := g^*\bigl(tI - J_n(t)\bigr)g,
    \qquad
    Q_t := Q_t^0 + (t + 1)e_+ + (t - 1)e_-.
    \label{eq:Qt-def}
\end{equation}
In the eigenbasis of $J_n(t)$, $Q_t$ is simply the weighted exponential sum $\sum_{z \in \mathbf{z}_t} (t - z) e_z$ (with $t - (-1) = t + 1$ and $t - 1$ as the two endpoint weights). Since the direction $v = g/\|g\|$ of a complex Gaussian vector is uniform on the simplex in squared moduli and independent of the Gamma-distributed radius $\|g\|^2$, standard Peano/simplex integration expresses the cumulative distribution function and sign expectation $\G(Q) := \E[\sgn Q]$ of $Q_t$ directly in terms of barycentric divided differences of $(y - t)_+^{m}$ on $\mathbf{z}_t$. Writing $\rho_t$ for the probability density of $Q_t$,~\eqref{eq:En-prime} becomes
\begin{equation}
    H_n'(t) = \G(Q_t) + 2(n - 1)t\,\rho_t(0) - n(1 - t^2)\,\rho_t'(0).
    \label{eq:Hprime-gauss}
\end{equation}

A tiny integration-by-parts identity now splits~\eqref{eq:Hprime-gauss} cleanly in two. Whenever $Y = X + ae$ with $a \neq 0$ and $e \sim \mathrm{Exp}(1)$ independent of $X$, one has $\rho_X = \rho_Y + a\rho_Y'$ and $\G(X) = \G(Y) - 2a\rho_Y(0)$. Applying this once with $a = t - 1$ (peeling off $(t-1)e_-$), once with $a = t + 1$ (peeling off $(t+1)e_+$), and once in sequence (peeling off both to reach $Q_t^0$) yields:

\begin{lemma}[Decomposition of $H_n'(t)$]
    \label{lem:decomp}
    For every integer $n \ge 1$ and every $t \in (-1,1)$,
    \begin{equation}
        H_n'(t)
        = \frac12\Bigl\{\G\bigl(Q_t^0 + (t+1)e_+\bigr) + \G\bigl(Q_t^0 + (t-1)e_-\bigr)\Bigr\}
          + \frac{n}{2}\,\Delta_n(t),
        \label{eq:comp}
    \end{equation}
    where
    \begin{equation}
        \Delta_n(t) := \G(Q_t) - \G(Q_t^0) = 4t\,\rho_t(0) - 2(1 - t^2)\,\rho_t'(0).
        \label{eq:delta-def}
    \end{equation}
\end{lemma}

The first term in~\eqref{eq:comp} falls immediately to an exact symmetry argument.

\begin{proposition}[Strict positivity of the first term]
    \label{prop:first-term}
    Let $D_n^0 := \diag(\frac12, 1, \dots, 1, \frac12)$, $A := g^* D_n^0 g$, $Z := g^* J_n(0) g$, and $e \sim \mathrm{Exp}(1)$ independent of $g$. Then for all $n \ge 1$ and all $t > 0$,
    \[
        \frac12\Bigl\{\G\bigl(Q_t^0 + (t+1)e_+\bigr) + \G\bigl(Q_t^0 + (t-1)e_-\bigr)\Bigr\}
        = \Prob\bigl(|e - Z| < t(A + e)\bigr) > 0.
    \]
\end{proposition}

\begin{proof}
Observe that $tI - J_n(t) = tD_n^0 - J_n(0)$, so $Q_t^0 = tA - Z$. Hence $Q_t^0 + (t+1)e_+$ has the same law as $(e - Z) + t(A + e)$. For the second summand, let $S := \diag(1, -1, 1, -1, \dots, (-1)^n)$. Because $J_n(0)$ is zero on the diagonal and supported on the first off-diagonals, we have $S J_n(0) S = -J_n(0)$ and $S D_n^0 S = D_n^0$. Since $Sg \overset{d}{=} g$, it follows that $(A, -Z) \overset{d}{=} (A, Z)$, and therefore
\[
    Q_t^0 + (t-1)e_-
    = tA - Z + (t-1)e_-
    \overset{d}{=} tA + Z + (t-1)e
    = (Z - e) + t(A + e).
\]
Set $X := t(A + e) > 0$ (almost surely) and $Y := e - Z$. Then the two terms have laws $X + Y$ and $X - Y$. Almost surely $|Y| \neq X$; when $|Y| < X$, both $X+Y$ and $X-Y$ are positive so $\sgn(X+Y)+\sgn(X-Y) = 2$, while when $|Y| > X$ they have opposite signs so the sum is $0$. Taking expectations yields $\frac12(\E\sgn(X+Y) + \E\sgn(X-Y)) = \Prob(|Y| < X)$.

To see strict positivity, evaluate at the explicit point $g_0 = (1, 1, 0, \dots, 0)^T \in \C^{n+1}$ and $e_0 = 1$: then $Z_0 = g_0^* J_n(0) g_0 = 1 = e_0$, so $|e_0 - Z_0| = 0 < t(A_0 + e_0)$. By continuity of $(g, e) \mapsto |e - Z| - t(A + e)$ and strict positivity of the joint Gaussian--exponential density on $\C^{n+1} \times (0, \infty)$, the open set $\{|e - Z| < t(A + e)\}$ has strictly positive probability.
\end{proof}

By Lemma~\ref{lem:decomp} and Proposition~\ref{prop:first-term}, proving $H_n'(t) > 0$ on $(0, 1/2]$ reduces entirely to conquering one last fortress: showing that $\Delta_n(t) > 0$.

\section{The Turning Point: Exact Characteristic Function and the Vertical-Chord Integral}
\label{sec:chord}

Earlier asymptotic approaches treated $Q_t$ as a generic $(n+3)$-dimensional quadratic form, combining trace expansions with local central limit theorems; near $t = 0$, that route hit a threshold around $n \sim 10^{40}$. Yet glancing back at Section~\ref{sec:matrix} reveals that $Q_t$ is far from generic: its characteristic function can be evaluated \emph{verbatim} as a closed rational function.

For $s = i\xi$ ($\xi \in \R \setminus \{0\}$), the Gaussian--exponential determinant formula gives
\[
    \Phi_{n,t}(\xi) := \E\bigl[e^{i\xi Q_t}\bigr]
    = \Bigl[\bigl(1 - s(t+1)\bigr)\bigl(1 - s(t-1)\bigr)\det\bigl(I_{n+1} - s(tI - J_n(t))\bigr)\Bigr]^{-1}.
\]
Factor out $(-s)^{n+3}$ and set $y := t - 1/s = t + i/\xi \in \C \setminus [-1,1]$; the denominator is precisely $(-s)^{n+3}(y^2 - 1)p_t(y)$. Now apply the classical \emph{Joukowsky map}
\[
    y = \frac{u + u^{-1}}{2}, \qquad u \in \D \setminus \{0\},
\]
and write $w := u^{-1}$. Substituting $U_m(y) = \frac{w^{m+1} - w^{-m-1}}{w - w^{-1}}$ into~\eqref{eq:charpoly}, the three terms miraculously collapse into a difference of two squares:
\[
    U_{n+1}(y) - 2tU_n(y) + t^2 U_{n-1}(y)
    = \frac{w^n(w - t)^2 - w^{-n}(w^{-1} - t)^2}{w - w^{-1}}.
\]
At the same time, $-s = (y - t)^{-1}$ yields
\begin{equation}
    s = \frac{-2u}{1 - 2tu + u^2},
    \qquad
    \frac{ds}{du} = \frac{-2(1 - u^2)}{(1 - 2tu + u^2)^2}.
    \label{eq:s-u}
\end{equation}
Multiplying all factors together, every power of $2$ and $u$ cancels cleanly to produce the exact characteristic function:
\begin{equation}
    \Phi_{n,t}(\xi)
    = \frac{(1 - 2tu + u^2)^{n+3}}{(1 - u^2)\bigl[(1 - tu)^2 - u^{2n+2}(u - t)^2\bigr]}.
    \label{eq:phi-exact}
\end{equation}
Pause for a moment to appreciate the anatomy of~\eqref{eq:phi-exact}: in the denominator, $(1 - tu)^2$ is the infinite-dimensional half-line Jacobi limit, whereas $-u^{2n+2}(u - t)^2$ is the finite-dimensional boundary reflection coming back from the far end of the matrix. Inside the unit disk $\D$, the difference between finite order $n$ and the $n \to \infty$ limit is \emph{exponentially small}.

Now insert~\eqref{eq:phi-exact} into the Fourier inversion formula for $\Delta_n(t)$:
\[
    \Delta_n(t) = 4t\,\rho_t(0) - 2(1 - t^2)\rho_t'(0)
    = \frac{1}{2\pi i}\int_{-i\infty}^{i\infty}\bigl(4t + 2(1 - t^2)s\bigr)\Phi_{n,t}(s/i)\,ds.
\]
(Even if $t$ happens to equal a node $u_j(t)$, the matrix $tI - J_n(t)$ has at most one zero eigenvalue because the eigenvalues are simple, so $\Phi_{n,t}(\xi) = O(|\xi|^{-(n+2)})$ and $\xi\Phi_{n,t} \in L^1(\R)$ for all $n \ge 1$.)

Changing variables via~\eqref{eq:s-u} triggers four favorable cancellations in rapid succession:
\begin{enumerate}[label=\textup{(\arabic*)}]
    \item The factor $(1 - u^2)$ in the denominator of $\Phi_{n,t}$ cancels identically against the numerator of $ds/du$.
    \item The linear factor simplifies to $4t + 2(1 - t^2)s = \frac{4(t - u)(1 - tu)}{1 - 2tu + u^2}$, introducing no pole at $u = 0$.
    \item Introducing the Blaschke disk automorphism $\varphi_t(u) := \frac{u - t}{1 - tu}$, we have $|\varphi_t(u)| < 1$ and $|u| < 1$ throughout $\D$, hence $|u^{2n+2}\varphi_t(u)^2| < 1$: the integrand is \emph{holomorphic with no poles} in the open unit disk $\D$. At the two boundary endpoints $u_\pm := t \pm ia$ ($a := \sqrt{1 - t^2}$), the denominator has only a simple zero on $\partial\D$, which is crushed by the $n$-th-order zero $(1 - 2tu + u^2)^n$ in the numerator, leaving bounded removable singularities.
    \item By Cauchy's theorem, we may deform the integration curve inside $\D$ into the straight \emph{vertical chord} $u = t + iv$ ($v \in [-a, a]$) joining $u_-$ and $u_+$. Along this chord,
    \[
        1 - 2t(t + iv) + (t + iv)^2 = a^2 - v^2 \ge 0:
    \]
    the imaginary part vanishes identically!
\end{enumerate}

Using conjugate symmetry across the real axis ($v \leftrightarrow -v$) yields the central analytic engine of the paper.

\begin{proposition}[Vertical-chord integral for $\Delta_n(t)$]
    \label{prop:chord}
    For every integer $n \ge 1$ and every $t \in (0,1)$, write $a := \sqrt{1 - t^2}$ and $\varphi_t(u) := \frac{u - t}{1 - tu}$. Then
    \begin{equation}
        \Delta_n(t)
        = \frac{8}{\pi}\int_0^a (a^2 - v^2)^n\,
          \RePart\!\left(\frac{-\varphi_t(u)}{1 - u^{2n+2}\varphi_t(u)^2}\right)_{u = t + iv} dv.
        \label{eq:delta-chord}
    \end{equation}
\end{proposition}

\begin{remark}[Direct contour-free proof in Lean~4]
    \label{rem:chord-lean}
    The formalization in \code{ChordIntegral.lean} proves~\eqref{eq:delta-chord} directly from the barycentric divided-difference formula for $\Delta_n(t)$ (via \code{barycentricDelta_eq_chordDelta}), bypassing Fourier inversion of probability densities altogether. Under the Joukowsky map, the vertical chord $u(v) = t + iv$ ($v \in [-a,a]$) traces a closed loop $\Gamma(v) = \frac12(u(v) + u(v)^{-1})$ winding clockwise once around $(t,1]$ and zero times around $[-1,t)$. Expanding the rational integrand in partial fractions over the $n+3$ extended nodes $\mathbf{z}_t$ reduces the entire identity to elementary single-pole integrals $\int \frac{\Gamma'(v)}{\Gamma(v) - z}\,dv$, each evaluated by an explicit \code{Complex.log} primitive and the Fundamental Theorem of Calculus.
\end{remark}

\section{Main Term versus Reflection Tail: All Orders \texorpdfstring{$n \ge 4$}{n >= 4} at Once}
\label{sec:n4}

Write $z := u^{2n+2}\varphi_t(u)^2$ for the reflection term in~\eqref{eq:delta-chord}, and split the integrand into a half-line main term and a reflection tail:
\[
    \frac{-\varphi_t(u)}{1 - z} = -\varphi_t(u) + \frac{-\varphi_t(u)\,z}{1 - z}.
\]
Along the vertical chord $u = t + iv$, abbreviate $D := |1 - tu|^2 = a^4 + t^2 v^2$ and $x := |u|^2 = t^2 + v^2$. Then $a^2 - v^2 = 1 - x$, and
\[
    -\varphi_t(u)
    = \frac{-iv}{a^2 - itv}
    = \frac{tv^2 - ia^2 v}{D}
    =: \alpha + i\beta.
\]

A quick order-of-magnitude check already reveals why the proof works effortlessly:
\begin{itemize}[leftmargin=1.8em]
    \item The main-term real part $\alpha = \frac{tv^2}{D} > 0$ is strictly positive everywhere; weighted by $(a^2 - v^2)^n$, its integral on $t \in (0, 1/2]$ scales like $(1 - t^2)^n \ge (3/4)^n$.
    \item For the tail, the algebraic identity
    \[
        x - \frac{v^2}{D}
        = x - \frac{x - t^2}{1 - (2 - x)t^2}
        = \frac{t^2(1 - x)^2}{D} \ge 0
    \]
    shows that $|\varphi_t(u)|^2 = \frac{v^2}{D} \le x$, hence $|z| \le x^{n+2}$. Weighted by $(1 - x)^n$, the product $(1 - x)^n x^{n+2}$ peaks at $4^{-n}$!
\end{itemize}

One subtlety must be respected when bounding the tail: as $t \to 0^+$, the main term scales like $O(t)$. A naive bound $|\varphi_t|\frac{|z|}{1 - |z|}$ discards the factor $t$, causing the tail-to-main ratio to blow up as $t \to 0^+$. In reality, at $t = 0$ both $u = iv$ and $\varphi_0(iv) = iv$ are purely imaginary, so $z = (iv)^{2n+4}$ is purely real and $-\varphi_0(iv)\frac{z}{1-z}$ is purely imaginary---its real part vanishes identically at $t = 0$!

To preserve this $O(t)$ cancellation, we take the real part \emph{before} applying the triangle inequality. Writing $\frac{z}{1 - z} = \frac{z - |z|^2}{|1 - z|^2}$ gives
\begin{equation}
    \left|\RePart\!\left((\alpha + i\beta)\frac{z}{1 - z}\right)\right|
    \le \alpha\,\frac{|z|}{1 - |z|} + |\beta|\,\frac{|\ImPart z|}{(1 - |z|)^2}.
    \label{eq:real-sep}
\end{equation}
The first term already carries $\alpha = \frac{tv^2}{D}$. For the second term, write $\arg u = \frac{\pi}{2} - \theta_1$ ($\theta_1 := \arctan\frac{t}{v}$) and $\arg \varphi_t(u) = \frac{\pi}{2} + \theta_2$ ($\theta_2 := \arctan\frac{tv}{a^2}$). Adding arguments turns all right angles into an integer multiple of $\pi$:
\[
    \arg z = (n + 2)\pi + \psi,
    \qquad
    \psi := 2\theta_2 - (2n + 2)\theta_1.
\]
Since $\theta_1, \theta_2 > 0$ have opposite signs in $\psi$, and since $\arctan y < y$ for $y > 0$ while $v \le a \implies \frac{v}{a^2} \le \frac{1}{v}$, we obtain
\[
    |\psi| \le \max\bigl(2\theta_2, (2n+2)\theta_1\bigr) < \frac{(2n + 2)t}{v}.
\]
Consequently $|\ImPart z| = |z|\,|\sin\psi| < \frac{(2n+2)t}{v}|z|$. Multiplying by $|\beta| = \frac{a^2 v}{D}$ cancels the factor $v$ and extracts the common factor $t$ cleanly!

Dividing~\eqref{eq:delta-chord} by $\frac{8t}{\pi}$ yields the uniform comparison
\begin{equation}
    \frac{\pi}{8t}\Delta_n(t) \ge M_n(t) - R_n(t)
    \qquad \bigl(t \in (0, 1/2]\bigr),
    \label{eq:Mn-Rn}
\end{equation}
where
\[
    M_n(t) := \int_0^a (a^2 - v^2)^n \frac{v^2}{D}\,dv,
    \qquad
    R_n(t) := \int_0^a (a^2 - v^2)^n\left[\frac{v^2}{D}\frac{|z|}{1 - |z|}
              + \frac{(2n+2)a^2}{D}\frac{|z|}{(1 - |z|)^2}\right]dv.
\]

\begin{proposition}[All orders $n \ge 4$]
    \label{prop:n4}
    For every integer $n \ge 4$ and every $t \in (0, 1/2]$, we have $\Delta_n(t) > 0$.
\end{proposition}

\begin{proof}
We bound $M_n(t)$ from below and $R_n(t)$ from above by closed-form Beta functions.

\smallskip
\noindent\textbf{Step 1: Closed-form lower bound on $M_n(t)$.}
Substituting $v = aw$ gives
\[
    M_n(t) = a^{2n+1}\int_0^1 \frac{(1 - w^2)^n w^2}{a^2 + t^2 w^2}\,dw.
\]
With respect to the probability weight proportional to $(1 - w^2)^n w^2\,dw$ on $[0,1]$, whose total mass is $\frac12 B(\frac32, n+1)$, the expectation of $w^2$ is $\frac{B(5/2, n+1)}{B(3/2, n+1)} = \frac{3}{2n+5}$. Applying Jensen's inequality (equivalently Cauchy--Schwarz) to the convex function $\lambda \mapsto \lambda^{-1}$ gives
\[
    M_n(t)
    \ge \frac{a^{2n+1}}{a^2 + \frac{3}{2n+5}t^2}\cdot\frac12 B\!\left(\frac32, n+1\right)
    = \frac{2n+5}{2\bigl((2n+5) - (2n+2)t^2\bigr)}\,(1 - t^2)^{n + 1/2}\,B\!\left(\frac32, n+1\right).
\]
One checks directly that $(1 - s)^{n + 1/2}/((2n+5) - (2n+2)s)$ is strictly decreasing in $s = t^2 \in [0, 1/4]$ (its logarithmic derivative has numerator $-(2n+1)(2n+5) + (2n+2)(2n-1)s + 4(n+1) < 0$), so the right-hand side attains its minimum at $t = 1/2$:
\begin{equation}
    M_n(t) \ge M_{\mathrm{lb}}(n)
    := \frac{2n + 5}{3(n + 3)}\left(\frac34\right)^{n + 1/2} B\!\left(\frac32, n + 1\right).
    \label{eq:Mlb}
\end{equation}

\smallskip
\noindent\textbf{Step 2: Closed-form upper bound on $R_n(t)$.}
In $R_n(t)$, change variables to $x = t^2 + v^2$ ($dv = \frac{dx}{2v}$, with $v \le \sqrt{x}$) and extend the lower integration limit from $t^2$ down to $0$. For all $t^2 \le \frac14$ and $x \in [0,1]$, elementary algebra gives
\[
    D = 1 - (2 - x)t^2 \ge \frac{2 + x}{4},
    \qquad
    \frac{a^2}{D^2} \le \frac{12}{(2 + x)^2},
    \qquad
    1 - |z| \ge 1 - x^{n+2} = (1 - x)S_{n+1}(x),
\]
where $S_{n+1}(x) := \sum_{k=0}^{n+1} x^k$. For every $n \ge 2$ and $x > 0$, two applications of AM--GM (using $3^2 = 9 > 8 = 2^3$, so $3^{2/3} > 2$) yield
\[
    (2 + x)S_{n+1}(x)
    \ge (2 + x)(1 + x + x^2 + x^3)
    = 3x + \Bigl(\frac43 + 3x^2\Bigr) + \Bigl(\frac13 + \frac13 + 3x^3\Bigr) + x^4
    > 9x.
\]
Using $|z| \le x^{n+1}\cdot\frac{v^2}{D}$ (so that in the first term $\frac{v^2}{D}|z| \le x^{n+2}\cdot\frac{v^2}{D}$, and in the second term $v \le \sqrt{x}$ combines with $|z| = x^{n+1}\frac{v^2}{D}$ to cancel $1/v$ in $dv = \frac{dx}{2v}$ cleanly: $\frac{a^2}{D}\frac{|z|}{(1-|z|)^2}\,dv = \frac{a^2}{2D^2}\frac{x^{n+1} v}{(1-|z|)^2}\,dx \le \frac{a^2}{2D^2}\frac{x^{n+3/2}}{(1-|z|)^2}\,dx$), substituting $D \ge \frac{2+x}{4}$, $\frac{a^2}{D^2} \le \frac{12}{(2+x)^2}$, and $(2+x)S_{n+1}(x) > 9x$ gives
\begin{equation}
    R_n(t) < T_1(n) + T_2(n)
    := \frac29\, B\!\left(n,\, n + \frac52\right)
     + \frac{4(n + 1)}{27}\, B\!\left(n - 1,\, n + \frac12\right).
    \label{eq:Rub}
\end{equation}

\smallskip
\noindent\textbf{Step 3: Base case $n = 4$ and geometric decay.}
Define the dimensionless ratios $A_k(n) := T_k(n)/M_{\mathrm{lb}}(n)$ ($k = 1, 2$); the factor $\sqrt{\pi}$ cancels identically. At $n = 4$, exact Beta evaluation gives
\[
    A_1(4) + A_2(4)
    = \frac{4039168\sqrt{3}}{9183213}
    \approx 0.76183 < 1
\]
(exact square $\frac{16314878132224}{28110467001123} < 1$, with strict margin $> 0.238$). By the Beta recurrence $B(p+1, q) = \frac{p}{p+q}B(p,q)$, the step ratios $r_k(n) := A_k(n+1)/A_k(n)$ are explicit rational functions of $n$ (displayed in Appendix~\ref{app:poly-n4}) whose shifted polynomials at $n = m + 4$ ($m \ge 0$) have strictly positive coefficients, proving $r_1(n) < \frac12$ and $r_2(n) < \frac12$ for all $n \ge 4$ (asymptotically both tend to $\frac13$). By induction on $n \ge 4$,
\[
    \frac{R_n(t)}{M_n(t)}
    \le 2^{-(n-4)}\bigl(A_1(4) + A_2(4)\bigr)
    < 0.76183 \cdot 2^{-(n-4)} < 1
    \qquad \bigl(t \in (0, 1/2]\bigr),
\]
whence $\Delta_n(t) > 0$.
\end{proof}

\section{Low Orders (\texorpdfstring{$n = 1, 2, 3$}{n = 1, 2, 3}) and Proof of Theorem~\ref{thm:full}}
\label{sec:finish}

Only the three lowest orders $n = 1, 2, 3$ remain.

\begin{proposition}[Low-order strict monotonicity]
    \label{prop:low-orders}
    For $n \in \{1, 2, 3\}$ and every $t \in (0,1)$, we have $E_n'(t) < 0$.
\end{proposition}

\begin{proof}
\textbf{Order $n = 1$.} The nodes of $J_1(t)$ are $u_0(t) = \frac{t-1}{2}$ and $u_1(t) = \frac{t+1}{2}$, giving $E_1(t) = \frac{1 - t^2}{2}$ and $E_1'(t) = -t < 0$ on $(0,1)$.

\smallskip
\noindent\textbf{Order $n = 2$.} The eigenvalues of $J_2(t)$ are $\frac{t - \sqrt{t^2+8}}{4}$, $\frac{t}{2}$, $\frac{t + \sqrt{t^2+8}}{4}$, which yields
\[
    6E_2(t) = (2 + 7t^2)\sqrt{8 + t^2} - 2 - 24t^2 - t^3,
    \qquad
    6E_2'(t) = \frac{3t}{\sqrt{8 + t^2}}\Bigl(7t^2 + 38 - (t + 16)\sqrt{8 + t^2}\Bigr).
\]
Comparing squares gives the polynomial identity
\[
    (t + 16)^2(8 + t^2) - (7t^2 + 38)^2
    = 4(t + 1)\bigl(151 - 87t + 20t^2 - 12t^3\bigr).
\]
On $[0,1]$, $151 - 87t + 20t^2 - 12t^3 = (151 - 99t) + 20t^2 + 12t(1 - t^2) \ge 52 > 0$, so $E_2'(t) < 0$ on $(0,1)$.

\smallskip
\noindent\textbf{Order $n = 3$.} Exterior monotonicity on $[1/2, 1)$ is already covered by Proposition~\ref{prop:outer}. On $(0, 1/2]$, we remark first that an elementary algebraic proof already gives the unique center maximum: setting $C_\pm := 3 \pm 13t + 17t^2 \pm 15t^3$ and $r_\pm := \sqrt{5 \pm 2t + t^2}$,~\eqref{eq:En} expands to
\[
    E_3(t) = \frac{C_+ r_+ + C_- r_-}{8} - 1 - 12t^2 - 2t^4,
\]
and square-root concavity / Taylor bounds yield $E_3(t) - E_3(0) \le -\frac{11}{32}t^2 < 0$ for all $t \in (0, 1/2]$ (formalized as \code{explicitError_orderedNode_three_lt_center}). To obtain the stronger strict monotonicity $E_3'(t) < 0$ on all of $(0,1)$, we split $(0, 1/2]$ into five subintervals $[t_0, t_1]$ and certify $\Delta_3(t) > 0$ on each block via a tightened one-dimensional Riemann-sum upper bound on $R_3(t)$ versus the local Jensen lower bound $M_{\mathrm{lb}}(3; t_1)$, as detailed in Appendix~\ref{app:n3}. Combining $\Delta_3(t) > 0$ with Lemma~\ref{lem:decomp} and Proposition~\ref{prop:first-term} gives $E_3'(t) < 0$ on $(0, 1/2]$ as well.
\end{proof}

\begin{proof}[Proof of Theorem~\ref{thm:full}]
(i) Combining Proposition~\ref{prop:outer} (exterior $[1/2,1)$), Lemma~\ref{lem:decomp} + Proposition~\ref{prop:first-term} + Proposition~\ref{prop:n4} (interior $(0,1/2]$ for $n \ge 4$), and Proposition~\ref{prop:low-orders} ($n = 1, 2, 3$), we have $E_n'(t) < 0$ for every $n \ge 1$ and every $t \in (0,1)$. Since $E_n$ is continuous on $[0,1]$, it is strictly decreasing on $[0,1]$. By evenness $E_n(-t) = E_n(t)$, $E_n(t) < E_n(0)$ for every $t \in [-1,1] \setminus \{0\}$, and $E_n(0)$ equals~\eqref{eq:bn-thm} by~\eqref{eq:bn}.

(ii) On $I = [-1,1]$, combining~\eqref{eq:Tn-bound},~\eqref{eq:peano}, and~(i) gives, for all $F \in W^{n+1,1}(I)$,
\begin{align*}
    \Lone{F^{(n)}}
    &\le \Lone{T_n F} + \Lone{F^{(n)} - T_n F} \\
    &\le A_n\Lone{F} + \int_{-1}^{1} E_n(t)\bigl|F^{(n+1)}(t)\bigr|\,dt
     \le 2^n n!\,\Lone{F} + B_n\Lone{F^{(n+1)}}.
\end{align*}
For an arbitrary interval $[a,b]$ with $L = b - a > 0$ and $f \in W^{n+1,1}([a,b])$, define $F(y) := f\!\left(\frac{a+b}{2} + \frac{L}{2}y\right)$ on $[-1,1]$. By the chain rule and change of variables,
\[
    \Lone{F^{(m)}} = \left(\frac{L}{2}\right)^{m-1}\lVert f^{(m)}\rVert_{L^1([a,b])}
    \qquad (0 \le m \le n+1).
\]
Substituting $m = 0, n, n+1$ into the unit-interval inequality and dividing by $(L/2)^{n-1}$ gives~\eqref{eq:full-ab}, since $2^n n!\,(L/2)^{-n} = 4^n n!\,L^{-n}$.

(iii) Optimality of both constants on $C^\infty([a,b])$ follows from Propositions~\ref{prop:An} and~\ref{prop:sharp} via the same affine rescaling between $[-1,1]$ and $[a,b]$.

(iv) By affine rescaling it suffices to work on $I = [-1,1]$. If $F = \lambda U_n$, then $F^{(n+1)} \equiv 0$ and equality holds by Proposition~\ref{prop:An}. Conversely, suppose $F \in C^{n+1}(I)$ achieves equality in~\eqref{eq:main} with $A_n = 2^n n!$ and $B_n = E_n(0)$. Then every step in
\[
    \Lone{F^{(n)}}
    \le \Lone{T_n F} + \int_{-1}^{1} E_n(t)\bigl|F^{(n+1)}(t)\bigr|\,dt
    \le A_n\Lone{F} + E_n(0)\Lone{F^{(n+1)}}
\]
must be an equality, forcing
\[
    \int_{-1}^{1} \bigl(E_n(0) - E_n(t)\bigr)\bigl|F^{(n+1)}(t)\bigr|\,dt = 0.
\]
The integrand is continuous and nonnegative on $[-1,1]$, and $E_n(0) - E_n(t) > 0$ for all $t \neq 0$ by~(i). Therefore $F^{(n+1)}(t) = 0$ for every $t \neq 0$, and by continuity $F^{(n+1)} \equiv 0$ on $[-1,1]$. Thus $F$ is a polynomial of degree at most $n$ attaining equality in Proposition~\ref{prop:An}, which forces $F = \lambda U_n$ for some $\lambda \in \R$.
\end{proof}

\section{Formal Verification in Lean 4}
\label{sec:lean}

Every step of the proof of Theorem~\ref{thm:full}---from Chebyshev moment orthogonality and the Jacobi moving-node system to the Gaussian decomposition, the vertical-chord integral, the $n \ge 4$ Beta bounds, the low-order proofs for $n = 1, 2, 3$, and the sharp Sobolev/smoothness conclusions---has been formally verified in Lean~4.32.2 with Mathlib inside the directory \code{Lean/SharpL1/Analysis/}.

\subsection{Master entry point and axiom audit}

The top-level unconditional theorem is
\begin{center}
    \code{SharpL1.Analysis.theorem_sharp_L1_interpolation_full}
\end{center}
in \code{Lean/SharpL1/Analysis/MainTheorem.lean}. Indexed by \code{k : Nat} (corresponding to order $n = k + 1 \ge 1$) and arbitrary real endpoints \code{a < b}, it bundles all four parts of Theorem~\ref{thm:full}:
\begin{enumerate}[label=\textup{(\roman*)}]
    \item \code{explicitErrorDerivCurve_neg_all_orders}, \code{explicitErrorCurve_strictAntiOn_Icc_zero_one}, \code{explicitErrorCurve_zero_eq_centerBestConstant}, and \code{explicitErrorCurve_unique_global_max_zero};
    \item \code{sharp_L1_interpolation_W1} on the one-dimensional Sobolev class \code{SobolevW1 (k + 2) a b f} (iterated absolutely continuous functions satisfying the Fundamental Theorem of Calculus up to order $n+1$ with integrable $(n+1)$-st derivative) and \code{sharp_L1_interpolation_contDiffOn} for \code{ContDiffOn R (k + 2) f (Icc a b)};
    \item \code{sharp_L1_interpolation_constants_optimal} over globally smooth functions \code{ContDiff R infty f};
    \item \code{sharp_L1_interpolation_contDiffOn_eq_iff} characterizing equality in $C^{n+1}([a,b])$.
\end{enumerate}

Verification is reproducible from \code{Lean/} via \code{lake build} followed by \code{lake env lean Audit.lean}. Across the entire library:
\begin{itemize}[leftmargin=1.8em]
    \item there are zero \code{sorry}s, zero \code{admit}s, zero \code{native_decide}s, and zero custom axioms;
    \item \code{Audit.lean} prints the axiom footprint of every main theorem and confirms that all declarations depend solely on the standard Lean axioms \code{[propext, Classical.choice, Quot.sound]} (with the pure-kernel \code{decide} certificates for $n = 3$ depending only on \code{[propext]}).
\end{itemize}

\subsection{Correspondence between manuscript and Lean 4 proofs}

While the mathematical statements match Theorem~\ref{thm:full} verbatim, four modules streamline analytic detours to stay clean and elementary in Mathlib:
\begin{itemize}[leftmargin=1.8em]
    \item \textbf{Divided-difference / CDF bridge (\code{DividedCDFBridge.lean}).} Rather than invoking Hermite--Genocchi simplex integration for $0 \le T_n h_t(x) \le 1$, the Lean proof establishes an exact algebraic identity between barycentric divided differences of truncated powers and the cumulative distribution function at $0$ of weighted exponential sums (\code{weightedCdfZero_eq_barycentric_truncatedPower}), yielding \code{movingOperator_normTruncatedPower_mem_Icc} and the eq.~\eqref{eq:comp} bridge \code{lawSignMean_Hcomp_of_nodes}.
    \item \textbf{Real-pair Gaussian realization (\code{FirstTermPos.lean}).} Complex Gaussian quadratic forms $g^* M g$ are realized via pairs of independent standard real Gaussian vectors $\frac12(x^T M x + y^T M y)$ (\code{pairQuadForm_pencil_law}), using orthogonal invariance of the standard real Gaussian measure and the explicit witness $x_0 = (1, 1, 0, \dots, 0)$, $y_0 = 0$, $e_0 = 1/2$ to prove \code{firstTerm_lawSignMean_pos}.
    \item \textbf{Direct partial-fraction chord integral (\code{ChordIntegral.lean}).} As noted in Remark~\ref{rem:chord-lean}, \code{barycentricDelta_eq_chordDelta} is proved directly via Lagrange partial fractions and elementary \code{Complex.log}/FTC integrals along the Joukowsky image of the vertical chord.
    \item \textbf{Kernel-checked fixed-point arithmetic for $n = 3$ (\code{LowOrderMono.lean}).} Instead of trusting external floating-point interval arithmetic (the companion \code{mpmath.iv} script \code{certify_chebyshev_L1.py} reported in Appendix~\ref{app:n3}), the Lean proof certifies the exact same five subintervals of $(0, 1/2]$ using natural-number fixed-point Riemann sums ($2^4, 2^4, 2^5, 2^5, 2^7$ subintervals, $10^9$ outer scale, $10^6$ intermediate scale, and $10^3$-scaled certified square-root bounds) evaluated inside the Lean kernel by \code{decide} (\code{n3_block0_certified} through \code{n3_block4_certified}), concluding \code{chordDelta_three_pos}.
\end{itemize}

Table~\ref{tab:lean-map} summarizes the module-by-module correspondence. We emphasize the exact scope boundary: all of Theorem~\ref{thm:full} and its complete proof (Sections~\ref{sec:poly}--\ref{sec:finish} and Appendix~\ref{app:tech}) are formally verified in Lean~4, whereas the structural extensions in Appendix~\ref{app:extensions} are expository complements not required for Theorem~\ref{thm:full} and are not formalized.

\begin{table}[htbp]
\centering
\caption{Mapping between paper results and verified Lean~4 declarations in \code{Lean/SharpL1/Analysis/}.}
\label{tab:lean-map}
\footnotesize
\renewcommand{\UrlFont}{\ttfamily\footnotesize}
\renewcommand{\arraystretch}{1.2}
\begin{tabular}{@{}p{3.4cm}>{\raggedright\arraybackslash}p{3.7cm}>{\raggedright\arraybackslash}p{8.0cm}@{}}
\toprule
\textbf{Paper result} & \textbf{Lean module} & \textbf{Key verified declarations} \\
\midrule
Prop.~\ref{prop:An} (first constant $A_n$ \& $U_n$ uniqueness) &
\code{ChebyshevFirstConstant.lean} &
\code{intervalIntegral_chebyshevSign_mul_pow_eq_zero}, \code{polynomial_L1_deriv_le}, \code{firstBestConstant_le_of_poly_bound}, \code{polynomial_L1_deriv_eq_iff} \\
\addlinespace[3pt]
Prop.~\ref{prop:jacobi} (global Jacobi nodes) &
\code{GlobalNodes.lean} (+ \code{Jacobi*.lean}) &
\code{jacobi_charpoly}, \code{hasDerivAt_orderedNode}, \code{orderedNode_bijOn_Icc}, \code{orderedNode_neg}, \code{orderedNode_zero} \\
\addlinespace[3pt]
Sec.~\ref{sec:sharp} (operator $T_n$, $E_n$, eq.~\eqref{eq:bn}) &
\code{MovingOperator.lean} &
\code{movingOperator_L1_le}, \code{movingOperator_poly}, \code{movingOperator_normTruncatedPower_mem_Icc}, \code{SobolevW1.deriv_sub_movingOperator_eq}, \code{integral_abs_indicator_sub_movingOperator_eq_explicitError}, \code{explicitError_center_eq} \\
\addlinespace[3pt]
Prop.~\ref{prop:sharp} \& $C^{n+1}$ equality &
\code{SharpSobolev.lean} &
\code{sharp_sobolev_Icc_le}, \code{sharp_firstConstant_Icc_le}, \code{sharp_secondConstant_Icc_le}, \code{sharp_sobolev_contDiffOn_Icc_eq_iff} \\
\addlinespace[3pt]
Prop.~\ref{prop:outer} \& Lem.~\ref{lem:decomp} &
\code{ExplicitError.lean}, \code{DividedCDFBridge.lean} &
\code{explicitErrorDerivCurve_neg_of_Ico_half_one}, \code{deriv_HCurve_eq_lawSignMean}, \code{lawSignMean_Hcomp_of_nodes} \\
\addlinespace[3pt]
Prop.~\ref{prop:first-term} (first term $> 0$) &
\code{FirstTermPos.lean} &
\code{pairQuadForm_pencil_law}, \code{firstTerm_lawSignMean_pos} \\
\addlinespace[3pt]
Prop.~\ref{prop:chord} (vertical-chord integral) &
\code{ChordIntegral.lean} &
\code{barycentricDelta_eq_chordDelta}, \code{lawSignMean_diff_eq_chordDelta} \\
\addlinespace[3pt]
Prop.~\ref{prop:n4} \& App.~\ref{app:poly-n4} ($n \ge 4$) &
\code{ChordBoundsN4.lean} &
\code{rem_lt_main_of_ge_four}, \code{chordDelta_pos_of_ge_four} \\
\addlinespace[3pt]
Prop.~\ref{prop:low-orders} \& App.~\ref{app:n3} ($n=1,2,3$) &
\code{LowOrderMono.lean} &
\code{explicitError_orderedNode_one_strictAntiOn}, \code{explicitError_orderedNode_two_strictAntiOn}, \code{explicitError_orderedNode_three_lt_center}, \code{n3_block0_certified}--\code{n3_block4_certified}, \code{chordDelta_three_pos} \\
\addlinespace[3pt]
\textbf{Theorem~\ref{thm:full}} (unconditional bundle) &
\code{MainTheorem.lean} &
\code{explicitErrorDerivCurve_neg_all_orders}, \code{explicitErrorCurve_unique_global_max_zero}, \code{sharp_L1_interpolation_W1}, \code{sharp_L1_interpolation_constants_optimal}, \code{sharp_L1_interpolation_contDiffOn_eq_iff}, \code{theorem_sharp_L1_interpolation_full} \\
\bottomrule
\end{tabular}
\end{table}

\section{Retrospective and Open Problems}
\label{sec:open}

Tracing the arc from the quadratic polynomial $x^2 - \frac14$ to the arbitrary-order Chebyshev polynomial $U_n$, the Jacobi matrix $J_n(x)$, the complex Gaussian quadratic form $Q_t$, and finally the Joukowsky vertical chord $u = t + iv$, one sees why the earlier asymptotic barrier near $n \sim 10^{40}$ dissolved so completely. Low orders did not harbor any hidden anomaly; rather, generic random-matrix bounds could not see the extraordinarily clean rational structure of $\Phi_{n,t}$ in Joukowsky coordinates. Once the contour is pulled onto the vertical chord $u = t + iv$, oscillation disappears, the $(3/4)^n$ versus $(1/4)^n$ gap surfaces immediately, the $t \to 0^+$ cancellation becomes transparent, and the complete proof turns out to be both shorter and far more elementary than a local central limit theorem.

Several stronger properties, though not needed for Theorem~\ref{thm:full}, remain inviting open questions:
\begin{enumerate}[label=\textup{(\arabic*)}]
    \item \textbf{Global strict concavity.} Is $E_n$ strictly concave on the entire open interval $(-1,1)$, i.e.\ $E_n''(t) < 0$ for all $t \in (-1,1)$? At present we know strict concavity at the origin ($E_n''(0) < 0$, Appendix~\ref{app:gaussian}) and asymptotically in the $t = O(n^{-1/2})$ scaling window.
    \item \textbf{Strict monotonicity of the best constants.} Is the sequence $B_n = E_n(0)$ strictly increasing in $n \ge 1$? The asymptotic expansion $1 - B_n \sim (\pi n)^{-1/2}$ (Appendix~\ref{app:asymptotics}) establishes monotonicity for sufficiently large $n$, but an all-order proof remains open.
    \item \textbf{Global positivity of $\Delta_n(t)$.} The density comparison inequality $\Delta_n(t) > 0$ (equivalently $\rho_t'(0)/\rho_t(0) < \frac{2t}{1-t^2}$ for the density $\rho_t$ of $Q_t$, or $-f_t'(t)/f_t(t) < \frac{2(n+1)t}{1-t^2}$ for the density $f_t$ of the uniform simplex convex combination of the $n+3$ extended nodes $\mathbf{z}_t$) is now proved for all $n \ge 3$ and $t \in (0, 1/2]$. Does it hold across the exterior $t \in (1/2, 1)$ and for $n = 1, 2$ as well?
\end{enumerate}

\appendix

\section{Expository Extensions: Generalized Gauss Quadrature, Center Asymptotics, and Local Gaussian Profile}
\label{app:extensions}

This appendix records three complementary structures surrounding the moving-node family. They are \emph{not} required for the proof of Theorem~\ref{thm:full} and are \emph{not} part of the Lean~4 formalization in Section~\ref{sec:lean}; we include them in expository form because they illuminate the broader landscape behind $J_n(x)$ and $E_n(t)$.

\subsection{Cross-order interlacing and generalized Gauss quadrature}

Setting $y = \cos\theta$ in~\eqref{eq:charpoly} and letting $\phi_x(\theta) := \arg\frac{e^{i\theta} - x}{1 - xe^{i\theta}}$ denote the Blaschke phase (which increases strictly from $0$ to $\pi$ on $[0,\pi]$), the eigenvalue angles of $u_{n,j}(x)$ satisfy the phase equation
\[
    (n + 1)\theta + \phi_x(\theta) = k\pi \qquad (k = 1, \dots, n+1).
\]
Increasing $n$ by $1$ adds $\theta \in (0, \pi)$ to the left-hand side, which forces strict interlacing of nodes across consecutive orders:
\[
    u_{n+1,j}(x) < u_{n,j}(x) < u_{n+1,j+1}(x) \qquad (j = 0, \dots, n).
\]
Furthermore, embedding the finite matrix $J_n(x)$ into the semi-infinite Jacobi operator on $\ell^2(\N)$ whose $(0,0)$-entry is $x/2$ and whose remaining diagonal entries vanish (with off-diagonals $\frac12$), any closed walk of length at most $2n$ starting at index $0$ never reaches the far boundary entry at index $n$. Hence the moments $\langle e_0, J_n(x)^m e_0 \rangle$ coincide with those of the half-line operator for all $0 \le m \le 2n$. Computing the Stieltjes transform of the half-line operator yields the \emph{generalized Gauss quadrature rule}
\[
    \sum_{j=0}^{n} u_j'(x)\,P\bigl(u_j(x)\bigr)
    = \int_{-1}^{1} P(y)\,\frac{2\sqrt{1 - y^2}}{\pi(1 - 2xy + x^2)}\,dy
    \qquad (\deg P \le 2n).
\]
At $x = 0$ the far-end perturbation disappears and precision increases by one ($\deg P \le 2n+1$), recovering classical Chebyshev--Gauss quadrature. Thus the moving nodes across all orders $n$ are successive interlacing discretizations of a single $x$-dependent spectral measure.

\subsection{Cancellation-free integral and asymptotics for \texorpdfstring{$B_n = E_n(0)$}{B\_n = E\_n(0)}}
\label{app:asymptotics}

Although~\eqref{eq:bn} gives $B_n$ in closed form, the prefactor $2^{n+1}$ causes severe numerical cancellation at large $n$. Applying Poisson summation to the alternating cosine sum together with Euler's reflection formula for the Gamma function yields the cancellation-free integral representation
\[
    1 - B_n
    = \frac{n + 2}{\pi(n + 1)}\int_0^1 \frac{s^{-1/2}(1 - s)^{n+1}}{1 - (-s)^{n+2}}\,ds.
\]
Since $s(1 - s) \le \frac14$ on $[0,1]$, replacing the denominator $1 - (-s)^{n+2}$ by $1$ introduces an error of order $O(4^{-n})$---precisely the $t = 0$ specialization of the $4^{-n}$ reflection tail encountered in Section~\ref{sec:n4}. Evaluating the leading Beta integral produces the high-precision asymptotic formula
\[
    1 - B_n
    = \frac{4^{n+2}}{\pi(n + 1)\binom{2n+4}{n+2}} + O(4^{-n})
    = \frac{1}{\sqrt{\pi n}}\left(1 + \frac{1}{8n} + O(n^{-2})\right).
\]
Table~\ref{tab:bn-values} lists exact formulas and numerical values for several orders $n$ (all now unconditionally proved to be the true optimal constants $B_n$).

\begin{table}[htbp]
\centering
\caption{Exact and numerical values of the optimal second constant $B_n = E_n(0)$.}
\label{tab:bn-values}
\begin{tabular}{rll@{\qquad}rll}
\toprule
$n$ & Exact $B_n$ & Numerical value & $n$ & Exact $B_n$ & Numerical value \\
\midrule
$1$ & $1/2$ & $0.500000000$ & $5$ & --- & $0.746758277$ \\
$2$ & $(2\sqrt{2} - 1)/3$ & $0.609475708$ & $10$ & --- & $0.820466444$ \\
$3$ & $3\sqrt{5}/4 - 1$ & $0.677050983$ & $20$ & --- & $0.873268528$ \\
$4$ & $(9\sqrt{3} - 12)/5$ & $0.717691454$ & $100$ & --- & $0.943514627$ \\
\bottomrule
\end{tabular}
\end{table}

\subsection{Local strict concavity at \texorpdfstring{$t = 0$}{t = 0} and the Gaussian scaling profile}
\label{app:gaussian}

A parallel Poisson-summation calculation shows that $E_n''(0) < 0$ for every $n \ge 1$, so the center $t = 0$ is always a strict local maximum with negative curvature. On the diffusive scale $t = r/\sqrt{n}$, the central limit theorem for $Q_{r/\sqrt{n}}/\sqrt{n} \Rightarrow N(r, \frac12)$ predicts a universal Gaussian scaling profile: uniformly on compact subsets of $r \in \R$,
\begin{align*}
    \sqrt{n}\,H_n(r/\sqrt{n}) &\longrightarrow r\erf(r) + \pi^{-1/2}e^{-r^2} = \E|r - Z|,\\
    H_n'(r/\sqrt{n}) &\longrightarrow \erf(r),\\
    n^{-1/2}H_n''(r/\sqrt{n}) &\longrightarrow 2\pi^{-1/2}e^{-r^2},
\end{align*}
where $Z \sim N(0, \frac12)$. The limit $\E|r - Z|$ gives an intuitive probabilistic explanation for the smooth parabolic minimum of $H_n$ at the origin.

\section{Technical Supplements: \texorpdfstring{$n \ge 4$}{n >= 4} Monotonicity Polynomials and \texorpdfstring{$n = 3$}{n = 3} Certification}
\label{app:tech}

To keep the proof of Theorem~\ref{thm:full} completely self-contained, we record the two verification details omitted from Sections~\ref{sec:n4} and~\ref{sec:finish}. Both are formally verified inside Lean~4 (\code{ChordBoundsN4.lean} and \code{LowOrderMono.lean}) and independently reproduced in seconds by the companion Python script \code{certify_chebyshev_L1.py}.

\subsection{Positive-coefficient polynomials for \texorpdfstring{$n \ge 4$}{n >= 4}}
\label{app:poly-n4}

From the Beta-function recurrence, the adjacent ratios $r_k(n) := A_k(n+1)/A_k(n)$ ($k = 1, 2$) in Proposition~\ref{prop:n4} evaluate explicitly to
\begin{align*}
    r_1(n) &= \frac{4n(n + 4)(2n + 5)^3}{3(n + 1)(n + 3)(2n + 7)(4n + 5)(4n + 7)},\\
    r_2(n) &= \frac{4(n + 2)(n - 1)(n + 4)(2n + 1)(2n + 5)^2}{3(n + 1)^2(n + 3)(2n + 7)(4n - 1)(4n + 1)}.
\end{align*}
Writing $n = m + 4$ ($m \ge 0$) and expanding $\mathrm{den}_k - 2\,\mathrm{num}_k$ gives
\begin{align*}
    \mathrm{den}_1 - 2\,\mathrm{num}_1
    &= 32m^5 + 912m^4 + 10354m^3 + 58855m^2 + 168954m + 198293 > 0,\\
    \mathrm{den}_2 - 2\,\mathrm{num}_2
    &= 32m^6 + 912m^5 + 10458m^4 + 61765m^3 + 199569m^2 + 342243m + 255933 > 0.
\end{align*}
Every coefficient is strictly positive, proving $r_1(n) < \frac12$ and $r_2(n) < \frac12$ for all $n \ge 4$.

\subsection{Five-subinterval certification of \texorpdfstring{$\Delta_3(t) > 0$}{Delta\_3(t) > 0} on \texorpdfstring{$(0, 1/2]$}{(0, 1/2]}}
\label{app:n3}

At $n = 3$, the global one-shot bound across all of $(0, 1/2]$ in Proposition~\ref{prop:n4} slightly exceeds $1$. Partitioning $(0, 1/2]$ into five compact blocks $[t_0, t_1]$ and writing $s_i := t_i^2$, we keep the exact lower integration limit $x \ge s_0$ and replace the global estimates by the sharper monotone block bounds
\[
    \frac{v^2}{D} \le \phi_{2,\max}(x) := \frac{x - s_0}{1 - (2 - x)s_0},
    \qquad
    1 - |z| \ge (1 - x)S_{\mathrm{eff}}(x),
\]
where
\[
    S_{\mathrm{eff}}(x) := 1 + x + x^2 + x^3 + x^4\,\frac{1 - s_0}{1 - (2 - x)s_0}.
\]
This produces an explicit smooth, bounded one-dimensional upper-bound integral $R_{\mathrm{ub}}(3; t_0, t_1)$ over the compact interval $[s_0, 1]$ (with zero truncation error and no tail) to compare against the local Jensen main-term lower bound
\[
    M_{\mathrm{lb}}(3; t_1)
    = \frac{176}{315}\,\frac{(1 - s_1)^{7/2}}{11 - 8s_1}.
\]

We verify $R_{\mathrm{ub}}(3; t_0, t_1) < M_{\mathrm{lb}}(3; t_1)$ on all five blocks in two completely independent ways:
\begin{enumerate}[label=\textup{(\arabic*)}]
    \item \textbf{Pure-kernel Lean~4 certification (\code{LowOrderMono.lean}).} On each block $[t_0, t_1]$, \code{LowOrderMono.lean} divides $[s_0, 1]$ into $2^d$ uniform subintervals ($d = 4, 4, 5, 5, 7$, i.e.\ $16, 16, 32, 32, 128$ boxes), bounds the integrand on each box by exact monotone natural-number fixed-point arithmetic (scale $10^9$ for box ceilings, $10^6$ for intermediate denominators, and $10^3$-scaled certified integer square-root enclosures), and checks the strict inequality against the integer lower bound on $M_{\mathrm{lb}}(3; t_1)$ inside the Lean kernel via \code{decide} (\code{n3_block0_certified} through \code{n3_block4_certified}, depending only on \code{[propext]}).
    \item \textbf{Companion Python interval arithmetic (\code{certify_chebyshev_L1.py}).} Using \code{mpmath.iv} at 40-digit precision with $128$ uniform subintervals per block produces the rigorous bounds reported in Table~\ref{tab:n3-cert}.
\end{enumerate}

\begin{table}[htbp]
\centering
\caption{Five-subinterval certification of $\Delta_3(t) > 0$ on $(0, 1/2]$ (Python \code{mpmath.iv} bounds and corresponding Lean~4 kernel \code{decide} grid sizes).}
\label{tab:n3-cert}
\footnotesize
\renewcommand{\UrlFont}{\ttfamily\footnotesize}
\setlength{\tabcolsep}{4.5pt}
\begin{tabular}{@{}ccccc@{}}
\toprule
Subinterval $[t_0, t_1]$ & $R_{\mathrm{ub}}$ bound & $M_{\mathrm{lb}}$ bound & Ratio (margin) & Lean grid ($2^d$) \\
\midrule
$[0,\, 1/4]$      & $0.018843402$ & $0.042453381$ & $\le 0.4438610$ ($>0.5561$) & $16$ (\code{n3_block0_certified}) \\
$[1/4,\, 7/20]$   & $0.019933669$ & $0.035293874$ & $\le 0.5647912$ ($>0.4352$) & $16$ (\code{n3_block1_certified}) \\
$[7/20,\, 21/50]$ & $0.020955782$ & $0.029542350$ & $\le 0.7093472$ ($>0.2906$) & $32$ (\code{n3_block2_certified}) \\
$[21/50,\, 23/50]$& $0.021203253$ & $0.026121386$ & $\le 0.8117200$ ($>0.1882$) & $32$ (\code{n3_block3_certified}) \\
$[23/50,\, 1/2]$  & $0.022212407$ & $0.022681618$ & $\le 0.9793132$ ($>0.0206$) & $128$ (\code{n3_block4_certified}) \\
\bottomrule
\end{tabular}
\end{table}

All five ratios are strictly less than $1$, establishing $\Delta_3(t) > 0$ for every $t \in (0, 1/2]$ and hence $E_3'(t) < 0$ for every $t \in (0, 1)$.

\end{document}
